\documentclass[11pt]{article}

\usepackage[margin=1in]{geometry}
\usepackage{microtype}
\usepackage{booktabs}
\usepackage{graphicx}
\usepackage{caption}
\usepackage{hyperref}
\hypersetup{
    colorlinks=true,
    urlcolor=blue!60!black,
    linkcolor=black,
    citecolor=black
}
\usepackage{parskip}
\usepackage{enumitem}
\setlist{itemsep=2pt,topsep=2pt}

% figures live alongside this .tex file
\graphicspath{{./}}

\title{Sub-Decision Prosodic Sensitivity in\\Qwen2-Audio-7B-Instruct\\
\large A diagnostic study using StressTest contrastive-stress items}
\author{Piyush Kumar}
\date{}

\begin{document}
\maketitle

\begin{abstract}
When a speech-LLM picks the wrong final answer on a contrastive-stress
question, two very different things might be happening: its audio
encoder might be discarding prosody (ears), or its language backbone
might be receiving the signal but underweighting it (brain). The
standard accuracy metric collapses both into a single ``wrong.'' This
study reads Qwen2-Audio-7B-Instruct's sub-decision logits --- the raw
probabilities behind its top-1 answer --- on the 218 items of the
StressTest benchmark, and shows that the model's internal
distribution \emph{does} shift in response to acoustic stress (paired
Wilcoxon $p \approx 10^{-18}$ in Format A, $10^{-17}$ in Format B),
even when the final answer is wrong. A text-rescue probe that feeds
the \emph{same language-model backbone} the transcript with the
stressed word in ALL CAPS lifts accuracy to $\sim$59\% --- beating
the primary audio model on both argmax (51--56\%) and pair-
consistency (14--20\%) metrics. The bottleneck is in the
audio$\rightarrow$LM interface, not in the language model. A
pre-registered ``stress median $>$ 95th percentile of strictest
floor'' check failed, but for a documented reason: Qwen2-Audio is
unusually fragile to $\pm 2\%$ time-stretch perturbations, so the
strictest floor is wider than the threshold assumed. We report this
transparently as a secondary finding about the model rather than
tune the bar to pass.
\end{abstract}

\section{Introduction}

Speech-LLMs are increasingly deployed in voice assistants,
transcription pipelines, and accessibility software. Prior work ---
most notably the StressTest benchmark (HUJI 2025) --- has shown that
these models often pick the wrong answer when meaning hinges on
which word the speaker stresses. But a wrong final answer is a
coarse signal. It could mean the model is truly deaf to prosody
(its internal probabilities don't move at all when stress changes),
or it could mean the model is partially sensitive (the probabilities
shift in the correct direction, just not enough to flip the top
choice). These two failure modes call for different engineering
responses: the first implicates the audio encoder, the second
implicates the language backbone or the audio$\rightarrow$LM
interface.

This project distinguishes them by reading the model's hidden
probability distribution instead of just its top-1 answer, and
complements that measurement with a text-rescue probe that isolates
the language model from the audio pathway. The whole study is
\textbf{inference-only} --- no training, no fine-tuning, a single
16~GB GPU.

The research question, in one sentence: \emph{when
Qwen2-Audio-7B-Instruct gets a contrastive-stress item wrong, does
the acoustic stress still shift the output distribution beneath the
final answer, and if so, where in the architecture is the signal
being lost?}

\section{Methods}

\subsection{Model and benchmark}

\begin{itemize}
    \item \textbf{Primary model.} Qwen2-Audio-7B-Instruct, loaded in
    4-bit NF4 quantization via bitsandbytes. Inference on a single
    Google Colab T4 (16~GB VRAM).
    \item \textbf{Language-only backbone.} Qwen2-7B-Instruct (same
    parameter count and tokenizer family as Qwen2-Audio's text
    decoder), used for the text-only, cascade, and text-rescue
    pathways.
    \item \textbf{Transcription model.} Whisper-large-v3 (cascade
    pathway only).
    \item \textbf{Benchmark.} StressTest (\texttt{slprl/StressTest}
    on HuggingFace), 218 audio items spanning 101 sentences with
    2--3 contrastive stress patterns each. Each item gives an audio
    clip, a gold transcript, a \texttt{possible\_answers} pair, a
    binary label, and a \texttt{stress\_pattern} dict marking the
    stressed word(s) by index. License CC-BY-NC-4.0, suitable for
    academic work.
\end{itemize}

\subsection{Logit-extraction strategy}

Rather than letting the model generate a free-form answer, we halt
after the \emph{very first} generated token and read the raw logits
for the two competing answer tokens. A two-way softmax over those
two logits gives a clean probability for each candidate. This
preserves the sub-decision distribution that argmax decoding
discards, and constrains the analysis to a well-defined
$P(A \mid A \text{ or } B)$ rather than a vocabulary-wide softmax.

Two answer formats were run to guard against single-format artifacts:

\begin{itemize}
    \item \textbf{Format A.} Options labeled \texttt{A.} /
    \texttt{B.}; next predicted token is \texttt{\ A} (id 362) or
    \texttt{\ B} (id 425).
    \item \textbf{Format B.} Options labeled \texttt{1.} /
    \texttt{2.}; next predicted token is \texttt{1} (id 16) or
    \texttt{2} (id 17). Format B's prompt ends with a trailing
    space so the space is absorbed into the prompt rather than the
    answer token (\texttt{\ 1} and \texttt{\ 2} are two tokens on
    this tokenizer).
\end{itemize}

Token IDs were empirically verified once on the Qwen2 tokenizer
(\texttt{scripts/verify\_tokens.py}) and hard-coded thereafter. An
independent \texttt{model.generate()} argmax-vs-extracted-logit-
argmax sanity check on a sampled subset prevents off-by-one position
errors and tokenizer drift.

\subsection{Pathways}

Five extraction pathways were run, each producing one JSONL record
per item per format (\texttt{results/logits\_*.jsonl}):

\begin{center}
\small
\begin{tabular}{p{2.4cm} p{6cm} p{6cm}}
\toprule
\textbf{Pathway} & \textbf{Input to the LM} & \textbf{Purpose} \\
\midrule
Primary & Qwen2-Audio sees raw audio & The thing being measured \\
Text-only & Qwen2-7B sees the gold transcript & Falsification: is the answer in the text? \\
Cascade & Whisper $\rightarrow$ Qwen2-7B & Falsification: does ASR + LM solve it? \\
Text-rescue & Qwen2-7B sees transcript with the stressed word in ALL CAPS & Architectural probe: ears vs.\ brain \\
Noise floor & Qwen2-Audio sees a meaning-preserving perturbation of the audio & Sub-decision baseline wobble \\
\bottomrule
\end{tabular}
\end{center}

The noise floor uses four perturbations applied to the original
audio: gain ($-3$~dB), additive Gaussian noise ($\sim$30~dB SNR),
time-stretch $\times 0.98$ (slow) and $\times 1.02$ (fast).
Perturbation seeds are deterministic per-item.

The text-rescue hint format is locked at
ALL-CAPS-on-stressed-word (e.g., \texttt{"I didn't say HE stole the
money."}). This convention matches the StressTest paper's own
usage and avoids leaking semantic framing through annotation
wording. Word indices come straight from the dataset's
\texttt{stress\_pattern["indices"]} field, so the capitalized word
is exactly the prosodically-stressed one and nothing else.

\subsection{Metrics and statistics}

Per item we compute the stress-induced shift
$|p_{\text{first}}(\text{item}) - p_{\text{first}}(\text{sister})|$,
averaged over sisters when a sentence has more than two stress
patterns. The corresponding noise-floor shift for perturbation $p$
is $|p_{\text{first}}(\text{clean}) - p_{\text{first}}
(\text{perturbed}_p)|$.

Two headline tests, one per answer format:

\begin{itemize}
    \item \textbf{Headline 1 --- decision-flip rate.} For pairs whose
    two members have \emph{different} gold labels ($n = 74$ eligible
    pairs), what fraction flips the model's argmax in the correct
    direction (each member lands on its own label)? Reported with
    bootstrap 95\% CIs over the indicator array.
    \item \textbf{Headline 2 --- paired Wilcoxon, stress shift vs.\
    strictest floor.} The ``strictest'' perturbation is the one
    with the largest 95th-percentile shift; we use a one-sided
    test of ``stress shift $>$ floor shift'' on the per-item
    paired values.
\end{itemize}

A pre-registered threshold check sits on top of Headline~2: the
median stress shift must exceed the 95th percentile of the
strictest floor's distribution. This is a strictly stronger bar
than the Wilcoxon and was committed to in \texttt{CLAUDE.md}
before the analysis was run.

All bootstrap CIs use 10\,000 resamples and a fixed seed
(\texttt{20260615}).

\section{Results}

\subsection{Pathway accuracies}

Argmax accuracy and pair-consistency accuracy (both items of a
contrastive pair correct):

\begin{center}
\small
\begin{tabular}{l c c c c}
\toprule
\textbf{Pathway} & \textbf{A argmax} & \textbf{B argmax} & \textbf{A pair-both} & \textbf{B pair-both} \\
\midrule
Primary (Qwen2-Audio) & 50.9\% & 55.5\% & 13.5\% & 19.5\% \\
Text-only (gold transcript) & 49.5\% & 48.2\% & 8.3\% & 8.3\% \\
Cascade (Whisper $\rightarrow$ LM) & 49.5\% & 48.2\% & 8.3\% & 8.3\% \\
\textbf{Text-rescue (ALL-CAPS hint)} & \textbf{59.6\%} & \textbf{59.2\%} & \textbf{23.3\%} & \textbf{22.6\%} \\
\bottomrule
\end{tabular}
\end{center}

Two falsification controls sit at chance: lexical content alone is
not enough to solve StressTest. Whisper transcribes these short
clean clips accurately enough that the cascade collapses to the
text-only condition (identical to three decimal places). The
primary audio model is a couple of points above chance on argmax
--- consistent with the StressTest paper --- and noticeably better
on pair-consistency, which can't be gamed by single-item heuristics
like option-order bias. The text-rescue row is the first piece of
the architectural-attribution story; we return to it in
Section~\ref{sec:attribution}.

\subsection{Headline 1 --- decision-flip rate}

Of the 74 within-sentence pairs whose two items carry
\emph{different} gold labels, how often does the model's argmax
flip in the correct direction?

\begin{center}
\small
\begin{tabular}{c c c}
\toprule
\textbf{Format} & \textbf{Flip rate} & \textbf{Bootstrap 95\% CI} \\
\midrule
A & 9.5\% & (4.1\%, 16.2\%) \\
B & 12.2\% & (5.4\%, 20.3\%) \\
\bottomrule
\end{tabular}
\end{center}

Both CIs exclude zero. The flip rate is modest but reliably
non-zero --- the model flips its answer in the correct direction
about one pair in ten. (A chance baseline would be a quarter ---
independent-coin-flip agreement on a pair of binary choices.)

\subsection{Headline 2 --- stress shift vs.\ noise-floor shift}

Paired one-sided Wilcoxon tests of ``stress shift $>$
strictest-floor shift'' on the 218 per-item paired observations:

\begin{center}
\small
\begin{tabular}{c l c c c c}
\toprule
\textbf{Format} & \textbf{Strictest floor} & $\boldsymbol{+/-/=}$ & \textbf{Median} $\boldsymbol{\Delta}$ & \textbf{95\% CI on median} $\boldsymbol{\Delta}$ & \textbf{Wilcoxon} $\boldsymbol{p}$ \\
\midrule
A & stretch\_fast & 163 / 55 / 0 & 0.082 & (0.031, 0.147) & $4.2 \times 10^{-18}$ \\
B & stretch\_slow & 158 / 60 / 0 & 0.034 & (0.009, 0.081) & $7.7 \times 10^{-17}$ \\
\bottomrule
\end{tabular}
\end{center}

In both formats, $\sim$72--75\% of items show a larger stress shift
than the strictest noise-floor shift. The paired-difference median
is positive in every bootstrap resample. The sub-decision logits
respond more to stress than to acoustic perturbation, per item,
with extremely high confidence.

\begin{figure}[h]
    \centering
    \includegraphics[width=\textwidth]{fig1_shift_distributions.png}
    \caption{Per-format histogram of the stress-shift distribution
    (filled blue) overlaid with the four noise-floor distributions
    (step outlines). The dashed vertical line marks the 95th
    percentile of the strictest perturbation --- this is the
    pre-registered threshold (see Section~\ref{sec:prereg}). Stress
    has substantial mass past 0.4 where the floors do not, but its
    median sits below the strictest-floor 95th percentile.}
    \label{fig:shift-dists}
\end{figure}

\begin{figure}[h]
    \centering
    \includegraphics[width=\textwidth]{fig2_stress_vs_floor_scatter.png}
    \caption{Per-item scatter of stress shift ($x$) against
    strictest noise-floor shift ($y$), with a $y=x$ diagonal.
    Points below the diagonal have stress $>$ floor; 163/218 =
    75\% of Format~A points and 158/218 = 72\% of Format~B points
    fall there. Color indicates whether the primary audio model
    got that item correct; correct and wrong points are roughly
    symmetrically distributed around the diagonal, which matches
    the project hypothesis --- stress moves the logits whether or
    not the final answer ends up correct.}
    \label{fig:scatter}
\end{figure}

\subsection{Decision-flip rate by stress position}

Breaking the flip rate down by average sentence-stress position
(initial / medial / final third of the sentence) gives a
suggestive but under-powered pattern.

\begin{figure}[h]
    \centering
    \includegraphics[width=\textwidth]{fig3_flip_by_position.png}
    \caption{Decision-flip rate (both items of a label-differing
    pair correct) binned by the average word-fraction position of
    the two members' stressed words. Initial-stress pairs flip at
    15\% / 23\% (A / B), medial at 10\% / 13\%, final at 5\% /
    5\%. Error bars are 95\% bootstrap CIs. With $n = 13 / 40 /
    21$ pairs the CIs are wide and overlap heavily; the downward
    trend is consistent in both formats but not load-bearing.}
    \label{fig:position}
\end{figure}

If the pattern survives a larger sample it would suggest the
model uses sentence-initial stress more than sentence-final
stress, consistent with attention-decay accounts of the audio
encoder. We do not stake the project's conclusions on this
figure.

\subsection{Architectural attribution: ears, not brain}
\label{sec:attribution}

The text-rescue probe re-runs the contrastive-stress task on the
\emph{same} language-model backbone used in the text-only and
cascade pathways (Qwen2-7B-Instruct), but with the stressed word
capitalized in the transcript. Comparing the four LM-backbone
conditions to the primary audio model:

\begin{center}
\small
\begin{tabular}{l c c c c}
\toprule
\textbf{Backbone / input} & \textbf{A argmax} & \textbf{B argmax} & \textbf{A pair-both} & \textbf{B pair-both} \\
\midrule
Qwen2-7B, gold transcript (no hint) & 49.5\% & 48.2\% & 8.3\% & 8.3\% \\
Qwen2-7B, Whisper transcript (no hint) & 49.5\% & 48.2\% & 8.3\% & 8.3\% \\
Qwen2-7B, gold transcript + ALL-CAPS hint & \textbf{59.6\%} & \textbf{59.2\%} & \textbf{23.3\%} & \textbf{22.6\%} \\
Qwen2-Audio, raw audio & 50.9\% & 55.5\% & 13.5\% & 19.5\% \\
\bottomrule
\end{tabular}
\end{center}

Three observations:

\begin{enumerate}
    \item \textbf{The LM backbone can use stress.} When stress
    arrives as a lexical cue (one word in caps), the same backbone
    that sat at chance on plain transcripts jumps $\sim$10 points
    on argmax and roughly doubles pair-consistency.
    \item \textbf{Text-rescue exceeds the primary audio model on
    both metrics, in both formats.} A trivial lexical encoding of
    stress is more useful to the downstream LM than whatever
    signal Qwen2-Audio currently extracts from the acoustic
    waveform.
    \item \textbf{The bottleneck is perceptual, not integrative.}
    The audio model is not deaf to stress --- the sub-decision
    Wilcoxon shows the logits respond --- but it is not routing
    that signal cleanly into the part of the representation the
    LM head reads from. ``Ears failed'' is shorthand for: stress
    reaches the audio embedding only weakly, and the bridge
    between audio encoder and language head loses most of it.
\end{enumerate}

This narrows the design space for any follow-up: improvements
should target the audio$\rightarrow$LM interface (e.g.\ richer
prosodic features routed to the text decoder, or training the
bridge with stress-contrastive objectives), not the language
model itself.

\subsection{The pre-registered threshold that failed}
\label{sec:prereg}

The pre-registration committed to a stricter bar than the
Wilcoxon: the median stress shift must exceed the 95th percentile
of the strictest floor. That check failed in both formats:

\begin{center}
\small
\begin{tabular}{c c l c}
\toprule
\textbf{Format} & \textbf{Stress median} & \textbf{Strictest 95th} & \textbf{Pass?} \\
\midrule
A & 0.20 & 0.38 (stretch\_fast) & FAIL \\
B & 0.10 & 0.31 (stretch\_slow) & FAIL \\
\bottomrule
\end{tabular}
\end{center}

The threshold was designed assuming all four perturbations would
behave like gain and additive noise --- small, tightly-distributed
shifts. They did not. The four noise-floor 95th percentiles are:

\begin{center}
\small
\begin{tabular}{l c c}
\toprule
\textbf{Perturbation} & \textbf{A} & \textbf{B} \\
\midrule
gain ($-3$~dB) & 0.034 & 0.029 \\
noise ($\sim$30~dB SNR) & 0.141 & 0.112 \\
stretch $\times 0.98$ & 0.326 & 0.313 \\
stretch $\times 1.02$ & 0.375 & 0.300 \\
\bottomrule
\end{tabular}
\end{center}

Time-stretch --- even at $\pm 2\%$, which is acoustically
inaudible as a meaning change --- moves the logits by an order of
magnitude more than gain or additive noise. An initial pilot used
$\pm 5\%$ stretches; tightening to $\pm 2\%$ only dropped
\texttt{stretch\_fast}'s 95th percentile from 0.475 to 0.375, and
\texttt{stretch\_slow} barely moved (0.322 $\rightarrow$ 0.326).
The model's internal representations are sensitive to durational
changes a listener would not flag.

This is itself a finding about Qwen2-Audio, not an artifact of
bad perturbation design. We report it transparently and do not
tune the threshold to clear it post hoc. The Wilcoxon and
decision-flip results stand on their own --- they are
\emph{per-item paired comparisons}, robust to whatever the floor
distribution looks like. The threshold test would have been a
useful additional check had it passed; its failure reflects a
property of the model.

A robustness re-run with the floor restricted to \texttt{gain} +
\texttt{noise} (post hoc, clearly flagged) would let Format~A
clear the threshold cleanly and leave Format~B borderline. We
mention this without promoting it: the study's primary claims
stand on the paired Wilcoxon and the text-rescue attribution, not
on the threshold check.

\section{Limitations}

\begin{itemize}
    \item \textbf{Modest $N$.} 218 items, 74 label-differing
    pairs for the decision-flip rate. The Wilcoxon is well-powered
    because it is paired and per-item, but the flip-rate CIs are
    wide ($\pm 6$ percentage points) and the position-breakdown is
    genuinely under-powered.
    \item \textbf{One primary model.} A planned attempt at
    Qwen2.5-Omni-3B was time-boxed at two days in Week~2 and
    dropped. The conclusions are specific to Qwen2-Audio-7B-
    Instruct, not ``speech-LLMs in general''.
    \item \textbf{English only}, StressTest speakers only.
    \item \textbf{Inference only.} The study makes no claim about
    whether training-time interventions could close the gap
    between sub-decision sensitivity and final-answer accuracy.
    \item \textbf{Threshold failure.} Reported transparently, with
    the model-property interpretation, but the pre-registered
    check did not pass; conclusions do not lean on it.
\end{itemize}

\section{Conclusion}

Reading the sub-decision logits of Qwen2-Audio-7B-Instruct on the
StressTest benchmark separates two failure modes that argmax
accuracy collapses. The model \emph{is} prosody-sensitive at the
distributional level (paired Wilcoxon $p \approx 10^{-18}$ /
$10^{-17}$; decision-flip rate CIs exclude zero in both answer
formats), but it does not weight that signal heavily enough to
flip the final answer most of the time. A text-rescue probe shows
that the same language-model backbone can use stress when it
arrives as a capitalized word --- and in fact uses it more
effectively than the audio model uses acoustic stress. The
bottleneck for contrastive-stress reasoning in this model is the
audio$\rightarrow$LM interface, not the language model.

A secondary, model-property finding: Qwen2-Audio's internal
representations are fragile to $\pm 2\%$ time-stretches that are
acoustically imperceptible as meaning changes, large enough to
invalidate the project's pre-registered floor-threshold check.

\section*{Reproducibility}

All code, intermediate outputs, and figures live in the project
repository at
\href{https://github.com/piyukr2/stresstest-subdecision-probe}
{\texttt{github.com/piyukr2/stresstest-subdecision-probe}}.
Specifically:

\begin{itemize}
    \item \texttt{scripts/verify\_tokens.py} --- token-ID
    verification (one-off; ids hard-coded into downstream scripts).
    \item \texttt{scripts/smoke\_test.py} --- 10-item pipeline
    smoke test with the generate-vs-forward sanity check.
    \item \texttt{scripts/extract\_logits.py} --- primary audio
    extraction (\texttt{results/logits\_primary\_\{A,B\}.jsonl}).
    \item \texttt{scripts/extract\_logits\_text.py} --- text-only
    pathway (\texttt{results/logits\_textonly\_\{A,B\}.jsonl}).
    \item \texttt{scripts/extract\_cascade.py} --- Whisper
    $\rightarrow$ LM cascade.
    \item \texttt{scripts/noise\_floor.py} --- four perturbations
    $\times$ two formats $\times$ 218 items = 8 JSONLs.
    \item \texttt{scripts/extract\_text\_rescue.py} --- ALL-CAPS
    text-rescue probe.
    \item \texttt{scripts/analysis.py} --- pathway summaries,
    headline statistics, bootstrap CIs. Writes
    \texttt{results/analysis\_summary.json}.
    \item \texttt{scripts/figures.py} --- generates the three
    figures in \texttt{figures/}.
\end{itemize}

Bootstrap seed is \texttt{20260615}. Noise-perturbation seeds are
deterministic per-item index. Token IDs were empirically verified
once on the Qwen2 tokenizer and hard-coded thereafter, eliminating
runtime-lookup as a failure mode.

\end{document}

